\documentclass[openany]{book}
\input{notes_white}
\usepackage{mathtools}
\usepackage{pgfplots}
\settitle{Open Problems in Computer Science - Skolem's Problem}
\begin{document}
\title{\Large{\underemphasise{Open Problems in Computer Science}} \\ \vspace{5pt}
\large{Skolem's Problem}}
\author{Eklavya Raman \\ \vspace{5pt} er24ms024@iiserkol.ac.in \\ \vspace{5pt}}
\date{\today}
\maketitle
\tableofcontents
\listoftheorems[
    show=mytheorem,
    title={List of Theorems}]


\chapter*{Introduction}
Skolem's Problem is a long-standing open problem in the field of theoretical computer science and number theory. It concerns the decidability of whether a given linear recurrence sequence contains a zero term. More formally, given a linear recurrence relation defined by integer coefficients, the problem asks whether there exists an integer $n$ such that the $n$-th term of the sequence is zero. This problem has significant implications in various areas, including automata theory, formal languages, and the theory of computation. Despite considerable efforts, Skolem's Problem remains unsolved in its general form, and it continues to be a subject of active research.

\chapter{Preliminaries}
We define some basic concepts and notations that will be used throughout the discussion of Skolem's Problem.
\section{Linear Recurrence Sequences}
\begin{mydefinition}{Linear Recurrence Sequence}
A linear recurrence sequence is a sequence of numbers $\{a_n\}$ defined by a linear recurrence relation of the form:
\[a_n = c_1 a_{n-1} + c_2 a_{n-2} + \ldots + c_k a_{n-k}\]
for $n \geq k$, where $c_1, c_2, \ldots, c_k$ are fixed coefficients and the initial terms $a_0, a_1, \ldots, a_{k-1}$ are given.
\end{mydefinition}
We will also need the concept of characteristic polynomials associated with linear recurrence sequences.
\begin{mydefinition}{Characteristic Polynomial}
The characteristic polynomial of a linear recurrence relation defined by coefficients $c_1, c_2, \ldots, c_k$ is given by:
\[P(x) = x^k - c_1 x^{k-1} - c_2 x^{k-2} - \ldots - c_k\]
\end{mydefinition}
\begin{mytheorem}[Roots of Characteristic Polynomial]
The general solution of a linear recurrence sequence can be expressed in terms of the roots of its characteristic polynomial. If the characteristic polynomial has distinct roots $r_1, r_2, \ldots, r_m$ with multiplicities $m_1, m_2, \ldots, m_m$, then the general term of the sequence can be written as:
\[a_n = \sum_{i=1}^{m} \left( P_i(n) r_i^n \right)\]
where $P_i(n)$ is a polynomial of degree $m_i - 1$.
\end{mytheorem}
\begin{proof}
The proof follows from the theory of linear difference equations. By substituting the general form of the solution into the recurrence relation and using the properties of polynomials and exponentials, we can derive the coefficients of the polynomials $P_i(n)$ based on the initial conditions of the sequence. The details involve solving a system of linear equations derived from the initial terms and the recurrence relation.
\end{proof}
\begin{example}
Consider the linear recurrence relation defined by:
\[a_n = 2a_{n-1} + 3a_{n-2}\]
with initial conditions $a_0 = 1$ and $a_1 = 2$. The characteristic polynomial is:
\[P(x) = x^2 - 2x - 3 = (x -3)(x + 1)\]
The roots are $r_1 = 3$ and $r_2 = -1$. The general solution is:
\[a_n = A \cdot 3^n + B \cdot (-1)^n\]
To find the constants $A$ and $B$, we use the initial conditions:
For $n=0$:
\[a_0 = A + B = 1\]
For $n=1$:
\[a_1 = 3A - B = 2\]
Solving this system of equations, we find $A = \frac{3}{4}$ and $B = \frac{1}{4}$. Thus, the explicit formula for the sequence is:
\[a_n = \frac{3}{4} \cdot 3^n + \frac{1}{4} \cdot (-1)^n\]
We can verify that this formula satisfies the recurrence relation and initial conditions as follows - 
For $n=2$:
\[a_2 = \frac{3}{4} \cdot 3^2 + \frac{1}{4} \cdot (-1)^2 = \frac{3}{4} \cdot 9 + \frac{1}{4} \cdot 1 = \frac{27}{4} + \frac{1}{4} = 7\]
For $n=3$:
\[a_3 = \frac{3}{4} \cdot 3^3 + \frac{1}{4} \cdot (-1)^3 = \frac{3}{4} \cdot 27 + \frac{1}{4} \cdot (-1) = \frac{81}{4} - \frac{1}{4} = 20\]
Thus, the sequence generated by the recurrence relation is $1, 2, 7, 20, \ldots$, which matches the values computed using the recurrence relation.
\end{example}
\begin{mytheorem}[Some properties of linear recurrence sequences]
Let $\{a_n\}$ be a linear recurrence sequence defined by a linear recurrence relation of order $k$ with integer coefficients. Then the following properties hold:
\begin{enumerate}
    \item The sequence $\{a_n\}$ is eventually periodic modulo any integer $m$. That is, there exist integers $N$ and $p$ such that for all $n \geq N$, $a_{n+p} \equiv a_n \mod m$.
    \item If the characteristic polynomial has a root that is a root of unity, then the sequence $\{a_n\}$ is periodic.
    \item The set of indices $n$ for which $a_n = 0$ is a union of finitely many arithmetic progressions along with a finite set.
\end{enumerate}
\end{mytheorem}
\begin{proof}
The proofs of these properties rely on the structure of linear recurrence sequences and their characteristic polynomials.
\begin{enumerate}
    \item The eventual periodicity modulo $m$ follows from the fact that there are only finitely many possible states for the vector of the last $k$ terms of the sequence modulo $m$. By the pigeonhole principle, the sequence must eventually repeat.
    \item If the characteristic polynomial has a root that is a root of unity, then the corresponding term in the general solution will be periodic, leading to the overall periodicity of the sequence.
    \item The structure of the zeros of linear recurrence sequences can be analyzed using techniques from algebraic number theory and Diophantine equations, leading to the conclusion about their distribution.
\end{enumerate}
\end{proof}
\section{Diophantine Equations}
\begin{mydefinition}{Diophantine Equation}
A Diophantine equation is a polynomial equation where the variables are intended to take integer values. A classic example is the equation $ax + by = c$, where $a$, $b$, and $c$ are given integers, and the solutions $(x, y)$ are sought in integers.
\end{mydefinition}
\begin{mytheorem}[Hilbert's Tenth Problem]
Hilbert's Tenth Problem asks for an algorithm to determine whether a given Diophantine equation has an integer solution. It was proven by Yuri Matiyasevich in 1970 that no such general algorithm exists, meaning that the problem is undecidable.
\end{mytheorem}
The proof of the undecidability of Hilbert's Tenth Problem involves showing that every recursively enumerable set can be represented as the set of integer solutions to some Diophantine equation. This is achieved by constructing specific Diophantine equations that encode the behavior of Turing machines, thereby linking the problem to the Halting Problem, which is known to be undecidable. The details of the construction are intricate and involve techniques from number theory and computability theory.For computability theory, read Roger Hartley's book "Theory of recursive functions and effective computability".
\end{document}